\documentclass[11pt,a4paper]{article}

\usepackage[T1]{fontenc}
\usepackage[utf8]{inputenc}
\usepackage{lmodern}
\usepackage[margin=2.4cm]{geometry}
\usepackage{graphicx}
\usepackage{booktabs}
\usepackage{amsmath}
\usepackage{amssymb}
\usepackage{caption}
\usepackage{float}
\usepackage[hidelinks]{hyperref}
\usepackage{xcolor}

\graphicspath{{figs/}}
\captionsetup{font=small,labelfont=bf,width=0.92\textwidth}

\newcommand{\detone}{P2\_1}
\newcommand{\dettwo}{P2\_2}
\newcommand{\uA}{\,\textmu A}

\title{Investigation of the intermittent efficiency loss of the P2 (BASKET)
pad-Micromegas detectors on the Saclay cosmic bench}
\author{Alexandra Kallitsopoulou}
\date{11 July 2026}

\begin{document}
\maketitle

\begin{abstract}
\noindent
During the long cosmic-ray efficiency runs of the P2 (BASKET) pad-Micromegas
prototypes taken between 30~June and 10~July~2026, detector~1 (\detone)
changed abruptly from a stable working state ($\varepsilon \simeq 80\%$) to a
quasi-dead state ($\varepsilon \simeq 5\%$), interrupted by short windows of
almost full efficiency during which the detector also sparks. Detector~2
(\dettwo) shows a different, much milder anomaly: a slow efficiency ramp
during the first $\sim$1.5--2~h of its run, after which it is stable.
This note documents the observations, the cross-checks performed
(mesh-current correlation, pulse-amplitude and hit-rate evolution, spatial
uniformity, HV scans before and after the failure, reference-tracker and
electronics sanity checks), and a discussion of the possible causes. All
evidence points to a global, bistable loss of amplification gain of \detone,
equivalent to more than $\sim$95~V of effective mesh voltage, that appeared
during an intense discharge episode in the night of 30~June\,/\,1~July and has
worsened since. Environmental gain modulation, DAQ effects and the reference
telescope are excluded. The most plausible cause is an intermittent
high-impedance element (marginal contact or discharge-induced defect) in the
mesh HV path downstream of the power-supply monitoring, and a set of targeted
hardware checks is proposed. The \dettwo{} turn-on transient is consistent
with a normal charging-up/stabilisation process and is not considered
pathological.
\end{abstract}

\section{Introduction}
\label{sec:intro}

The P2 (BASKET) detectors are large fan-shaped pad-Micromegas prototypes read
out on 1280 pads each. Two units, \detone{} and \dettwo, are under test on the
Saclay cosmic bench, where the M3 reference telescope provides independent
track reconstruction so that the absolute detection efficiency of the device
under test can be measured as a function of position, time and operating
voltage.

A series of long ($\geq$10~h) overnight runs was taken to establish the
working point and the stability of the detectors. The first long run of
\detone{} (30~June) showed a stable plateau efficiency of $\sim$80\%. From
1~July onwards, however, \detone{} has been essentially inefficient
($\varepsilon_{\text{any}}\simeq 5\%$), with striking transient
\emph{recovery windows} of one to a few hours during which the efficiency
returns to almost its healthy value --- in the run of 7~July, exactly the
period in which the mesh HV channel was sparking. \dettwo, operated for the
first time on 9~July, reaches a stable 77\% but only after a slow
$\sim$1.5--2~h ramp from zero.

The purpose of this note is to reconstruct the failure timeline, to test the
observable consequences of the candidate explanations against the recorded
data (efficiency time series, CAEN HV monitoring, raw pad hits, HV scans),
and to converge on a ranked list of hypotheses and follow-up actions.
Section~\ref{sec:setup} summarises the detectors and the measurement;
Section~\ref{sec:obs} presents the observations; Section~\ref{sec:analysis}
the dedicated correlation analyses; Section~\ref{sec:discussion} the
interpretation; and Section~\ref{sec:conclusions} the conclusions and
recommended actions.

\section{Detector and measurement description}
\label{sec:setup}

\subsection{The P2 (BASKET) detectors}

P2 is a \emph{metallic, non-resistive} bulk-Micromegas detector: the anode is
segmented into pads (no resistive layer, no strips), arranged in a fan
(``basket'') geometry covering roughly $70\lesssim r \lesssim 600$~mm with
1280 pads routed through ten MEC8 connectors. The amplification mesh is held
by an insulation mask with $\sim$11.7k small ($\sim$0.5~mm) support pillars
on a $\sim$4~mm grid. The chamber is operated with Ar/iC$_4$H$_{10}$ (95/5)
at atmospheric pressure. Two HV channels per detector are supplied by a CAEN
mainframe: the mesh (amplification) voltage and the drift (cathode) voltage
(\detone: channels 1:0 and 1:1; \dettwo: 1:2 and 1:3). The CAEN monitor
values (vmon, imon) are logged continuously to \texttt{hv\_monitor.csv} at
0.5--2~s sampling.

On \detone{} connectors 1 and 10 are physically disconnected (dead); on
\dettwo{} connector 1 is disconnected and connectors 8--10 sit on a front-end
unit excluded from the run. The dead connectors are removed from the pad map
in all analyses, leaving 1024 live pads on \detone{}; quoted efficiencies
refer to the remaining \emph{active area} only. The pads are read out by
DREAM-based front-end units (FEUs); pad pulse amplitudes are recorded in ADC
counts together with a global trigger timestamp.

\subsection{Efficiency measurement}

The M3 telescope provides reference tracks; a good ray is required to have
$N_{\text{clus}}\geq 3$ and $\chi^2 < 5$ in both coordinates. For each good
ray extrapolated to the P2 plane inside the active area, two efficiency
estimators are formed:
\begin{itemize}
  \item \textbf{within $R$}: a P2 pad hit is found within the match radius
        of the extrapolated impact point (20~mm for \detone, 40~mm for
        \dettwo, chosen on the efficiency-vs-radius plateau);
  \item \textbf{has any}: the event contains at least one P2 pad hit
        anywhere (an upper bound, insensitive to residual alignment).
\end{itemize}
Efficiencies are computed in 30-min bins to obtain the time evolution.
Periods during which the mesh HV channel draws
$\text{imon}\geq 2$\uA{} are treated as sparks and vetoed ($-2$/$+10$~s
guard window), and events with $\geq$20 simultaneous pads (internal
discharges) are removed. The numbers quoted here are from the
\emph{un-vetoed} products unless stated otherwise, so that the behaviour
during spark periods remains visible.

\begin{figure}[H]
  \centering
  \includegraphics[width=0.49\textwidth]{effmap_det1_0630.png}\hfill
  \includegraphics[width=0.49\textwidth]{effmap_det2_0709.png}
  \caption{Detectors in their working state, for reference. Left:
  \detone{} pad-level efficiency map (``within 20~mm'') from the healthy run
  of 30~June. Right: \dettwo{} efficiency map (``within 40~mm'') from the run
  of 9~July (plateau part). The fan-shaped active area and the excluded
  connectors are visible.}
  \label{fig:effmaps}
\end{figure}

\section{Observations}
\label{sec:obs}

\subsection{Run overview}

Table~\ref{tab:runs} summarises the five long runs analysed here. The same
analysis chain (M3 alignment, efficiency maps, HV-spark QA, amplitude QA) was
run for every dataset.

\begin{table}[H]
  \centering\small
  \caption{Long runs used in this investigation. Efficiencies are medians of
  the 30-min bins over the active area (``any'' / ``within $R$'');
  the peak column gives the highest single bin. Spark counts are merged
  discharge events with mesh imon $\geq 2$\uA{} (stage-07 QA).}
  \label{tab:runs}
  \begin{tabular}{llllllll}
    \toprule
    Run (start) & Det. & Mesh & Dur. & Median eff. & Peak eff. & Sparks &
    Recovery window \\
    \midrule
    6-30 19:24 & \detone & 420 V & 14 h & 80.7 / 69.8 \% & 84.6 / 73.9 \% & 58 &
    up to 04:15, then dead \\
    7-3 17:30 & \detone & 440 V & 14 h & 5.6 / 2.2 \% & 44.6 / 20.1 \% & 104 &
    01:00--02:30 \\
    7-7 22:52 & \detone & 440 V & 10 h & 6.3 / 4.2 \% & 82.7 / 62.6 \% & 11 &
    02:52--04:52 \\
    7-9 19:11 & \detone & 430 V & 10 h & 3.6 / 0.5 \% & 4.6 / 0.8 \% & 1 &
    none \\
    7-9 19:11 & \dettwo & 430 V & 10 h & 77.1 / 74.4 \% & 78.5 / 75.5 \% & 7 &
    (slow turn-on 0--1.8 h) \\
    \bottomrule
  \end{tabular}
\end{table}

\subsection{Timeline of the \detone{} failure}
\label{sec:timeline}

\begin{description}
  \item[30 June, 19:24 -- 1 July, ~04:10.] \detone{} operates stably at mesh
  420~V with $\varepsilon_{\text{any}}\simeq 80\%$,
  $\varepsilon_{20\,\text{mm}}\simeq 70\%$
  (Fig.~\ref{fig:overlay_0630}, top). The mesh channel sparks at a moderate
  rate throughout. At $t\simeq4.1$~h (23:30) the channel undergoes an
  exceptional \emph{sustained} discharge: 270~s long, 9.9\uA{} peak,
  1795~\textmu C integrated charge --- two orders of magnitude above a
  typical spark ($\sim$5--15~\textmu C). The detector survives it with no
  visible efficiency change.
  \item[1 July, 04:10--05:30.] A dense series of discharges begins at
  $t\simeq8.96$~h ($\sim$04:22) and the efficiency collapses from 80\% to
  $<$5\% within two bins (Fig.~\ref{fig:overlay_0630}). The detector never
  recovers during the remaining $\sim$5~h of the run, while continuing to
  spark heavily.
  \item[2 July, 12:04 -- evening.] A 16-point mesh-HV scan
  (345--420~V) is taken. The detector behaves \emph{normally} during the
  whole scan, with a clean turn-on curve reaching
  $\varepsilon_{\text{reco}}=68\%$, $\varepsilon_{\text{any}}=82\%$ at 420 V
  (Fig.~\ref{fig:hvscans}, left). The HV was powered off at the end of the
  scan. This day proves the failure was, at that point, fully reversible.
  \item[3 July, 17:30.] Long run at an increased working point, mesh 440~V.
  Immediately after the ramp the mesh channel enters a violent discharge
  storm --- 104 spark events within the first 1.6~h, including a 114~s
  sustained one --- after which it goes quiet. The detector is inefficient
  \emph{from the very first bin} (5\%), storm included. One recovery window
  occurs at 01:00--02:30 (peak 44.6\%/20.1\%) with \emph{zero} sparks in it
  (Fig.~\ref{fig:overlay_0704}).
  \item[7 July, 22:52.] Long run, mesh 440~V. Dead at $\sim$5--7\% all
  night except a two-hour window 02:52--04:52 that peaks at 82.7\%/62.6\% ---
  i.e.\ the healthy efficiency of 30~June. \emph{All 11 sparks of the run
  fall inside this window} (Fig.~\ref{fig:overlay_0707}).
  \item[9 July, 19:11.] Combined \detone+\dettwo{} run at mesh 430~V with an
  appended HV scan. \detone{} is dead the whole night (median 3.6\%/0.5\%,
  single spark) and dead at \emph{every} point of the 395--430~V scan
  (Fig.~\ref{fig:hvscans}, right). \dettwo, in the same run and gas, ramps
  from zero to a stable 77\% within $\sim$1.8~h and stays there
  (Fig.~\ref{fig:overlay_det2}).
\end{description}

\begin{figure}[H]
  \centering
  \includegraphics[width=0.9\textwidth]{overlay_det1_long.pdf}
  \caption{\detone, run of 30~June (mesh 420~V). Top: efficiency vs time in
  30-min bins (orange: any hit; green: within 20~mm); pink vertical lines
  mark mesh discharges (imon $\geq2$\uA). Bottom: raw mesh current (thin)
  and its 30-min median (points). The efficiency is stable at 80\% for
  8.8~h and collapses at $\sim$04:15 in coincidence with the onset of a
  dense discharge series; it never recovers.}
  \label{fig:overlay_0630}
\end{figure}

\begin{figure}[H]
  \centering
  \includegraphics[width=0.9\textwidth]{overlay_det1_long2.pdf}
  \caption{\detone, run of 3--4~July (mesh 440~V). The discharge storm at the
  start (0--1.6~h, 104 events) follows the ramp to the new 440~V working
  point. The detector is inefficient from the first bin. The recovery window
  at $t\simeq7.5$--9~h (01:00--02:30) contains no spark at all.}
  \label{fig:overlay_0704}
\end{figure}

\begin{figure}[H]
  \centering
  \includegraphics[width=0.9\textwidth]{overlay_det1_long3.pdf}
  \caption{\detone, run of 7--8~July (mesh 440~V). The two-hour recovery
  window (02:52--04:52) reaches the healthy efficiency of 30~June
  (82.7\%/62.6\%), and \emph{all 11 sparks of the run occur inside it}. The
  raw mesh current also shows a visibly denser sub-threshold activity band
  during the window.}
  \label{fig:overlay_0707}
\end{figure}

\begin{figure}[H]
  \centering
  \includegraphics[width=0.9\textwidth]{overlay_det2_long1.pdf}
  \caption{\dettwo, run of 9--10~July (mesh 430~V). The mesh was at full
  voltage from the start of the run (vmon constant); the efficiency
  nevertheless needs $\sim$1.8~h to ramp from zero to its stable plateau of
  77\%. Note the generally higher sub-threshold mesh-current activity of
  this detector.}
  \label{fig:overlay_det2}
\end{figure}

\section{Analysis}
\label{sec:analysis}

\subsection{Are the sparks the cause or a symptom?}
\label{sec:sparkcorr}

The naive reading of the 7~July run --- ``the detector is only efficient
when it sparks'' --- suggests the discharges might cause the recovery (e.g.\
by burning off a short). The 3~July run breaks this: its recovery window
contains 0 of the run's 104 sparks, and conversely its violent spark storm
(first 1.6~h) coincides with a fully \emph{inefficient} period. The
consistent picture across all runs is the opposite causality:
\begin{quote}
\emph{sparking requires amplification gain. When the gain is restored the
detector both detects cosmics and (at 440~V, close to the spark limit)
discharges occasionally; when the gain is absent it can do neither.}
\end{quote}
The sparks in the 7~July window are therefore a \emph{symptom} of the
restored gain, not its cause. Supporting this, the mean sub-threshold mesh
current in that window is $73$~nA against $22$~nA outside it --- exactly the
extra current expected from real gas amplification of the cosmic flux ---
while the 30-min \emph{median} imon stays at the few-nA noise floor in both
states, i.e.\ there is no ohmic leakage path appearing or disappearing
(Sec.~\ref{sec:discussion}).

\subsection{Gain evidence: amplitude, rate and spectra}
\label{sec:gain}

If the low-efficiency state were a tracking, alignment or matching artefact,
the pad hits themselves would look normal. They do not.
Figure~\ref{fig:amp_0707} shows the median pad-hit amplitude and the pad-hit
rate against time for the 7~July run: inside the recovery window the median
amplitude rises from 175 to 562~ADC ($\times$3.2) and the hit rate from
0.24~Hz to up to 4~Hz ($\times$10, matching the M3-triggered cosmic rate).
The amplitude \emph{spectrum} inside the window (Fig.~\ref{fig:spec_0707})
develops the extended Landau-like tail up to ADC saturation that
characterises real minimum-ionising signals, whereas outside the window the
spectrum is a soft, steeply-falling distribution consistent with residual
noise just above threshold. The same comparison for the 30~June run
(before/after the 04:15 collapse) gives 397 vs 268~ADC and a rate drop from
3.7~Hz to 0.24~Hz. This is direct evidence that the inefficiency is a
\textbf{collapse of the gas amplification gain}, not a reconstruction
artefact.

\begin{figure}[H]
  \centering
  \includegraphics[width=0.88\textwidth]{amplitude_det1_long3.pdf}
  \caption{\detone, 7~July run: median (and 90th-percentile) pad-hit
  amplitude (top) and pad-hit rate (bottom) in 30-min bins. Both jump by a
  large factor exactly inside the efficiency window (shaded).}
  \label{fig:amp_0707}
\end{figure}

\begin{figure}[H]
  \centering
  \includegraphics[width=0.62\textwidth]{amp_spectra_det1_long3.pdf}
  \caption{\detone, 7~July run: normalised pad-hit amplitude spectra inside
  (orange) and outside (grey) the recovery window. Inside the window a hard
  signal tail extending to ADC saturation appears; outside, only a soft
  noise-like population remains.}
  \label{fig:spec_0707}
\end{figure}

\subsection{The recovery is global, not local}
\label{sec:spatial}

A failure of one HV sector, one connector or one front-end would produce a
spatially structured effect. Figure~\ref{fig:hitmap} compares the pad hitmap
inside and outside the 7~July window: inside the 2.25~h window \emph{all
1024 live pads} fire, with the smooth radial illumination profile of the
healthy detector; outside it the (10$\times$ smaller) hit population is
spatially structureless noise. The recovery switches the \emph{entire}
active area on and off coherently, pointing to a global parameter --- in a
Micromegas, effectively the mesh--anode field.

\begin{figure}[H]
  \centering
  \includegraphics[width=\textwidth]{hitmap_window_det1_long3.pdf}
  \caption{\detone, 7~July run: pad hitmap inside (left) and outside (right)
  the two-hour recovery window (log colour scale, grey = pad never fired).
  The whole fan lights up uniformly inside the window: the gain loss and
  recovery are global.}
  \label{fig:hitmap}
\end{figure}

\subsection{HV response before and after: an effective voltage deficit
$>$95~V}
\label{sec:hvscan}

The two mesh-HV scans bracket the sick state
(Fig.~\ref{fig:hvscans}). On 2~July the turn-on is textbook-like: the
efficiency rises smoothly from 19.6\% (any) at 345~V to 82.2\% at 420~V. On
9~July the same detector gives 1--3.4\% at \emph{every} voltage between
395 and 430~V --- lower than the healthy detector delivered at 345~V. In the
dead state the operating point at 430--440~V therefore lies \emph{more than
95~V below} the healthy turn-on curve in effective amplification voltage.
This scale immediately rules out environmental gain modulation:
temperature/pressure variations of the gas density modulate the effective
mesh voltage at the level of a few volts, not a hundred.

\begin{figure}[H]
  \centering
  \includegraphics[width=0.49\textwidth]{hvscan_det1_0702.png}\hfill
  \includegraphics[width=0.49\textwidth]{hvscan_det1_0709.png}
  \caption{\detone{} mesh-HV scans one week apart. Left: 2~July, healthy
  turn-on reaching 68\%/82\% (reco/any) at 420~V. Right: 9~July, the
  efficiency is $<$3.5\% at every point of the 395--430~V scan (note the
  vertical scale): the detector no longer responds to the applied mesh
  voltage.}
  \label{fig:hvscans}
\end{figure}

\subsection{Sanity checks: telescope, drift HV, electronics}
\label{sec:sanity}

\begin{itemize}
  \item \textbf{M3 reference.} The number of good reference rays per 30-min
  bin is stable throughout the 7~July run (3004--3235, mean 3121): the
  denominator of the efficiency, the trigger and the telescope
  reconstruction are steady while the P2 efficiency swings by a factor 15.
  \item \textbf{Drift field.} The drift channel of \detone{} (600~V) is
  featureless during all runs: vmon within 1~V of setpoint, imon compatible
  with zero (99th percentile $<$1~nA). The anomaly involves the mesh
  (amplification) stage only.
  \item \textbf{Front-end electronics and thresholds.} The pedestal QA of
  the runs shows no anomaly, hits continue to be recorded in the dead state
  (at low rate and soft amplitudes), and the DREAM FEUs also serve the M3
  chambers, which measure normally. A dead-time, threshold or
  synchronisation problem cannot produce a state in which recorded pulses
  \emph{exist} but have collectively lost their amplitude tail.
  \item \textbf{CAEN monitoring.} During dead periods the mesh vmon sits at
  its setpoint (439.8~V median on 7~July) with no excursions outside spark
  transients, and the baseline imon is identical in the dead and live states
  (medians 3--6~nA). Whatever removes the amplification field is invisible
  to the power-supply monitoring, i.e.\ it sits \emph{downstream} of it, or
  it does not draw current.
\end{itemize}

\subsection{\dettwo: a different, benign transient}
\label{sec:det2}

\dettwo{} was switched on at full mesh voltage at run start, yet its
efficiency needs $\sim$1.8~h to reach the plateau
(Fig.~\ref{fig:overlay_det2}). During the ramp the pad-hit \emph{rate}
grows from 1.5 to 3.2~Hz while the median amplitude of recorded hits is
already at its plateau value (640 vs 668~ADC) --- the signature of a slowly
rising gain moving a threshold-biased population into acceptance, classic
charging-up/stabilisation behaviour (insulating pillar charging and gas
flushing after closing the fresh chamber). After the ramp the detector is
flat at 77.1\% (any) for 8~h, sparking rarely (7 events). Apart from one
28~s, 162~\textmu C sustained discharge at 23:43, which produced no
after-effect, nothing anomalous is seen. \dettwo{} therefore does
\emph{not} share the \detone{} pathology; its slow start should be handled
operationally (start physics 2~h after HV-on, or pre-bias the chamber).

Importantly, \dettwo{} running normally \emph{in the same run, the same gas
bottle and the same DAQ} while \detone{} is dead excludes every common-mode
explanation for the \detone{} state (gas mixture at the supply, trigger,
DAQ, telescope, environment), provided the two chambers share the gas
distribution; a line-specific flow or leak problem on the \detone{} branch
remains formally possible and is cheap to check.

\section{Discussion}
\label{sec:discussion}

The observations that any explanation must accommodate are:
\begin{enumerate}
  \item the loss is a \textbf{global gain collapse} equivalent to
        $>$95~V of mesh voltage (Secs.~\ref{sec:gain}--\ref{sec:hvscan});
  \item it appeared \textbf{abruptly at 04:15 on 1~July during a dense
        discharge series}, one week into otherwise stable operation, hours
        after an exceptional 1795~\textmu C sustained discharge
        (Sec.~\ref{sec:timeline});
  \item it was \textbf{fully reversible on 2~July} (healthy 10~h HV scan),
        returned after the HV-off/440~V-ramp cycle of 3~July (with a
        discharge storm at the ramp), and has since worsened: recovery
        windows of 1.5--2~h on 3 and 7~July at 440~V, none on 9~July at
        430~V;
  \item recoveries reach \textbf{full healthy performance} within
        $\sim$1~h and decay on a similar timescale --- a bistable system
        with slow transitions, not a gradual drift;
  \item in the dead state the mesh channel \textbf{draws no anomalous
        current} ($\leq$ few nA baseline, unchanged), and vmon is nominal at
        the CAEN;
  \item sparking correlates with the \emph{efficient} state whenever the
        detector is efficient at 440~V, but recovery can occur without any
        spark (3~July);
  \item \dettwo{} is healthy in identical conditions.
\end{enumerate}

\noindent The candidate explanations rank as follows.

\paragraph{H1 --- Environmental modulation (T, p, humidity): excluded as
primary cause.} Gas-density gain modulation amounts to a few volts of
effective mesh voltage; here $>$95~V is missing (Sec.~\ref{sec:hvscan}).
The night-time clustering of the two recovery windows (01:00--02:30,
02:52--04:52) and of the original failure (04:15) is suggestive of an
environmental \emph{trigger} of the state transitions (e.g.\ lab temperature
acting on a marginal mechanical contact), but the 9~July run covers the same
hours with no window, so even as a trigger it is not reproducible.
Temperature/pressure logging on the bench would settle this cheaply.

\paragraph{H2 --- Gas supply problem on the \detone{} line: unlikely,
testable.} Heavy air or water contamination can kill the gain by tens of
percent, and a blocked/leaking line specific to \detone{} would spare
\dettwo. However: the death was instantaneous (minutes) in coincidence with
a discharge series, not a slow drift as a gas volume exchanges; the 2~July
recovery was equally step-like; and full performance returns within an hour
in the recovery windows, faster than a $\sim$l/h-flushed chamber of this
volume can renew its gas. A flow-meter check and, if available, a swap of the
gas lines between the two detectors would remove this hypothesis entirely.

\paragraph{H3 --- DAQ, thresholds, reference tracking: excluded.}
Section~\ref{sec:sanity}: the telescope rate is constant, pedestals are
normal, hits are recorded in both states, and the amplitude spectrum
transformation (loss of the MIP tail) is a gas-gain signature that no
digital effect reproduces.

\paragraph{H4 --- CAEN channel or cabling up to the monitoring point:
excluded by the data, swap to be safe.} vmon and imon are nominal in the
dead state. A fault \emph{inside} the CAEN downstream of its ADCs cannot be
fully excluded from data alone; swapping the \detone{} and \dettwo{} mesh
channels at the mainframe is a five-minute test that would move the problem
with the channel if it lived there.

\paragraph{H5 --- Discharge-induced fault in the mesh HV path or in the
amplification gap: the leading hypothesis.} Everything points to the
effective mesh potential at the amplification gap being absent while the
supply delivers nominal voltage and sees no load. Two variants are
compatible with the data:
\begin{itemize}
  \item \textbf{(a) Intermittent open/high-impedance contact} in the mesh HV
  chain downstream of the CAEN monitoring --- HV connector, filter box,
  feedthrough solder joint, or the internal spring/wire contact to the mesh
  frame. An open mesh floats and discharges through leakage; the gap field
  collapses globally with \emph{no} current signature (obs.~1, 5). Occasional
  reconnection (thermally or mechanically driven, plausibly at night ---
  obs.~4 and the night windows) restores full field within the RC/leakage
  timescale, and at 440~V the freshly restored gap sparks (obs.~6). The
  initial failure and the 3~July relapse both coincide with violent
  discharge activity (obs.~2, 3), which is exactly what damages a marginal
  joint --- and what a re-arcing joint produces. The progressive worsening
  (windows shrinking to zero) fits a degrading contact.
  \item \textbf{(b) Discharge-deposited defect in the gap} (debris or a
  carbonised pillar spot after the 1795~\textmu C event) that locally
  connects mesh and anode and charges up: harder to reconcile with the
  \emph{absence} of any extra mesh current in the dead state ($\leq$nA,
  identical to the live state) and with the perfectly global character of
  the loss, so it is disfavoured relative to (a), but a low-current
  bistable path cannot be excluded without opening or measuring the chamber.
\end{itemize}

The recovery-window efficiency reaching exactly the healthy plateau (82.7\%
vs 84.6\% peak bins) shows the detector is \emph{intact} when the field is
present: gas, mesh, pads and electronics are all healthy. The defect is in
delivering or holding the field.

\section{Conclusions and recommended actions}
\label{sec:conclusions}

\detone{} has, since 04:15 on 1~July~2026, been in a bistable state in which
the gas amplification gain is absent most of the time (efficiency
$\simeq$5\%, effective mesh-voltage deficit $>$95~V) and occasionally returns
to full healthy performance for 1.5--2~h; the frequency and duration of the
recoveries have decreased from run to run and were zero in the last run. The
failure and its one relapse both coincided with intense discharge episodes
(including a 270~s, 1.8~mC sustained discharge), while the recovery windows
show that every subsystem --- gas, mesh, pads, electronics, telescope ---
still performs nominally when the amplification field is present. The
inefficiency is invisible to the CAEN monitoring: the fault sits downstream
of it or draws no current. The leading hypothesis is an intermittent
high-impedance element in the mesh HV path (damaged contact or
discharge-induced defect); gas, environment, DAQ and telescope are excluded
as primary causes. \dettwo{} shows only a normal $\sim$2~h charging-up
transient and is fit for physics.

\noindent Recommended actions, in order of increasing intrusiveness:
\begin{enumerate}
  \item \textbf{Live two-state diagnostic.} The pad-hit rate (0.25~Hz dead
  vs 3--4~Hz alive) is an unambiguous real-time indicator of the state. Run
  it together with the 0.5~s HV monitor during all interventions below, so
  that any manipulation that flips the state is caught immediately.
  \item \textbf{Swap the mesh HV channels} of \detone{} and \dettwo{} at the
  CAEN: excludes the mainframe channel and the first cable segment.
  \item \textbf{Mechanical/thermal stimulation test:} with HV on and the
  live diagnostic running, gently manipulate (wiggle, warm, cool) the HV
  cable, connectors, filter box and feedthrough of \detone. A state flip
  during manipulation localises the marginal contact.
  \item \textbf{Electrical measurement with HV off:} capacitance and
  insulation measurement between the mesh line and the pad plane / drift
  from the HV connector. An open mesh line shows as a missing
  mesh--anode capacitance; a damaged joint may show as an unstable reading
  under slight mechanical load.
  \item \textbf{Check the \detone{} gas branch} (flow meter reading,
  bubbler) and, if trivial, swap the gas lines of the two detectors:
  closes hypothesis H2 definitively.
  \item \textbf{Add temperature/pressure/humidity logging} to
  \texttt{hv\_monitor.csv}: tests the night-time-trigger conjecture and is
  generally useful for gain normalisation.
  \item \textbf{After any repair}, re-establish the working point with an HV
  scan and keep the mesh at $\leq$430~V initially: the 440~V ramp of 3~July
  triggered the worst discharge storm in the dataset and preceded the
  degradation of the recovery windows. Consider soft-trip current training
  before physics running.
  \item If all external checks are negative, the fault is internal
  (feedthrough-to-mesh bond or in-gap defect) and the chamber has to be
  opened for inspection --- last resort.
\end{enumerate}

\vspace{2ex}
\noindent\textit{All figures and the numerical summaries in this note were
produced from the standard P2 QA pipeline products (stages 06/07) and two
dedicated correlation scripts; figure sources live in
\texttt{Analysis/reports/p2\_efficiency\_intermittency/}.}

\end{document}
